\section{A coordinate change that moves the normal forms}
\label{sec:deciding-coordinate-change}

Work first on a product of three complex discs with coordinates $(x_1,x_2,z)$.
The two divisors are $D_1=\{x_1=0\}$ and $D_2=\{x_2=0\}$.
Write $\lambda_i=dx_i/x_i$.
A logarithmic residue removes a normal differential after moving it to the first position:
\[
r_1(\lambda_1\wedge\lambda_2)=\lambda_2|_{D_1},\qquad
r_2(\lambda_1\wedge\lambda_2)=-\lambda_1|_{D_2}.
\]
Consequently the two successive residues of this form are $1$ and $-1$.
The normal order determines the sign.

Change coordinates by
\begin{equation}
\label{eq:two-normal-unit}
y_1=e^z x_1,\qquad y_2=x_2,\qquad z'=z.
\end{equation}
This is an invertible holomorphic change, with the same two divisors.
Its normal logarithmic form is
\[
\frac{dy_1}{y_1}\wedge\frac{dy_2}{y_2}
=\lambda_1\wedge\lambda_2+dz\wedge\lambda_2.
\]
The second term has one normal logarithmic factor.
It cannot equal $f(z)\lambda_1\wedge\lambda_2$ for a holomorphic function $f$.
Indeed, on the complement of the divisors,
$\lambda_1\wedge\lambda_2$ and $dz\wedge\lambda_2$ are distinct exterior basis elements.
Thus the geometric coordinate change does not preserve the chosen space of top normal forms.

There are two questions.
The first asks whether the full logarithmic residue complex is independent of coordinates.
The second asks how the smaller normal complex compares with that coordinate-independent complex.
The distinction persists even though the normal complex includes every collision stratum.
We will construct both comparisons and examine which homotopies respect the number of collisions.

The starting language is linear algebra, holomorphic functions, and differential forms.
All complexes are complexes of complex vector spaces or sheaves of such spaces.
A differential has degree $+1$.
Its degree-$q$ cohomology is the kernel of the differential in degree $q$ modulo its image from degree $q-1$.
For a complex $V$, the shift is $V[a]^j=V^{j+a}$, with differential $(-1)^a d_V$.
The wedge of forms of degrees $p,q$ changes by $(-1)^{pq}$ when their order is exchanged.
Tensor products obey the same sign rule.
All edge sets, sums over strata, and covers used below are finite.
No completed tensor product, support condition, or analytic limiting procedure is used.

\section{The coordinate-independent residue complex}
\label{sec:residue-sheaf}

Let $Y$ be a complex manifold and $D=\bigcup_{e\in E}D_e$ a labeled simple normal crossing divisor.
This means that its components are smooth and, locally, are distinct coordinate hyperplanes.
For $S\subseteq E$, put
$D_S=\bigcap_{e\in S}D_e$, with $D_\varnothing=Y$.
The remaining components cut out a divisor $D^S$ on $D_S$.
An empty intersection contributes the zero complex.

On a chart of $D_S$ with remaining normal coordinates $x_e$ and tangential coordinates $z_a$,
the sheaf $\Omega^1_{D_S}(\log D^S)$ is the free holomorphic module with basis
$d\log x_e=dx_e/x_e$ and $dz_a$.
Its $q$th exterior power is $\Omega^q_{D_S}(\log D^S)$.
The exterior derivative extends the usual derivative on functions and satisfies
$d(d\log x_e)=0$.
A sheaf assigns sections to each open set, with restrictions and unique gluing of compatible local sections.
The displayed local modules define such a sheaf because changes between adapted charts preserve logarithmic forms.
The next calculation proves this preservation and determines the residue maps.

For $e\notin S$, expand $\alpha=d\log x_e\wedge\beta+\gamma$ in the coordinate exterior basis,
with neither $\beta$ nor $\gamma$ containing $d\log x_e$.
Define the left-normal residue
\begin{equation}
\label{eq:residue-definition}
r_{S,e}:\Omega^q_{D_S}(\log D^S)
\longrightarrow\Omega^{q-1}_{D_{S\cup\{e\}}}(\log D^{S\cup\{e\}}),
\qquad r_{S,e}\alpha=\beta|_{D_{S\cup\{e\}}}.
\end{equation}
This coefficient is independent of the presentation.
In a coordinate exterior basis, terms containing $dx_e$ with a regular coefficient have zero residue,
because $dx_e=x_e\,d\log x_e$.
Every other ambiguity changes $\beta$ by a form whose restriction is zero.

An adapted coordinate change between charts has normal equations
\begin{equation}
\label{eq:adapted-change}
y_e=u_e x_{\pi(e)},\qquad u_e\in\mathcal O_Y^\times.
\end{equation}
Here $\pi$ is a permutation of the locally present components, constant on each connected overlap.
The tangential coordinates can also change.
The superscript $\times$ means that $u_e$ is nowhere zero.
Both coordinate systems are assumed to be charts, so the full coordinate map has an invertible Jacobian.
Their logarithmic differentials satisfy
\begin{equation}
\label{eq:unit-differential}
d\log y_e=d\log x_{\pi(e)}+\alpha_e,
\qquad \alpha_e=u_e^{-1}du_e,\qquad d\alpha_e=0.
\end{equation}
The form $\alpha_e$ is regular and holomorphic.
Its closedness follows from
$d(u_e^{-1}du_e)=-u_e^{-2}du_e\wedge du_e=0$.
The inverse coordinate change has the same form, so the logarithmic modules agree.

Write $i_S:D_S\hookrightarrow Y$ for inclusion.
The direct image sheaf $i_{S,*}\mathcal F$ has sections
$\mathcal F(V\cap D_S)$ on an open set $V\subseteq Y$.
Set
\begin{equation}
\label{eq:full-residue-complex}
\mathcal R_Y=
\bigoplus_{S\subseteq E}i_{S,*}\Omega^\bullet_{D_S}(\log D^S)[-2|S|],
\qquad D=d+R,\qquad R=\sum_{S,e\notin S}r_{S,e}.
\end{equation}
A $q$-form on $D_S$ has total degree $q+2|S|$.
Thus a residue has degree one, although it lowers the form degree by one.
The shifts are even, so $d$ is the usual exterior derivative on each stratum.

\begin{proposition}[Residue naturality and strict descent]
\label{prop:strict-residue-descent}
The operator $D$ satisfies $D^2=0$.
Every adapted coordinate change induces a degree-zero chain isomorphism on~\eqref{eq:full-residue-complex} by ordinary pullback of forms and relabeling of strata.
These isomorphisms compose strictly.
Consequently the local descriptions glue to the displayed complex of sheaves on $Y$.
\end{proposition}

\begin{proof}
The coefficient of $d\log x_e\wedge d\log x_f$ changes sign when the two factors are exchanged.
Restriction of its remaining coefficient to $D_{S\cup\{e,f\}}$ is independent of order.
Hence, for distinct $e,f\notin S$,
\[
r_{S\cup\{e\},f}r_{S,e}
=-r_{S\cup\{f\},e}r_{S,f}.
\]
Once $e$ has been removed, no second $e$-residue is defined.
There is therefore no repeated-edge contribution to $R^2$.
All other terms cancel in the displayed pairs.

For $\alpha=d\log x_e\wedge\beta+\gamma$, the derivative is
$d\alpha=-d\log x_e\wedge d\beta+d\gamma$.
The form $d\gamma$ still has zero residue along $D_e$.
Thus $r_{S,e}d=-dr_{S,e}$.
Together with $d^2=0$ and $R^2=0$, this proves $D^2=0$.

Consider a coordinate map $g$ from an $x$-chart to a $y$-chart with equations~\eqref{eq:adapted-change}.
It maps the source stratum indexed by $\pi(S)$ to the target stratum indexed by $S$.
Denote this restriction by $g_S$.
On a term $d\log y_e\wedge\beta+\gamma$, pullback gives
$(d\log x_{\pi(e)}+\alpha_e)\wedge g^*\beta+g^*\gamma$.
The two regular terms have zero residue along $x_{\pi(e)}=0$.
Therefore
\begin{equation}
\label{eq:pullback-residue}
r_{\pi(S),\pi(e)}g_S^*=g_{S\cup\{e\}}^*r_{S,e}.
\end{equation}
Pullback also commutes with $d$ by the chain rule on functions and the exterior product rule.
It preserves total degree because $|\pi(S)|=|S|$.
Thus it intertwines $D$.
The inverse coordinate map supplies its inverse.
The identities $(g\circ h)^*=h^*g^*$ and $\pi_{g\circ h}=\pi_h\pi_g$
show strict compatibility on triple overlaps.
Unique gluing of differential forms proves the sheaf assertion.
\end{proof}

The permutation signs in this proposition occur inside wedge pullback.
For two normal directions, swapping $y_1=x_2$, $y_2=x_1$ sends
$d\log y_1\wedge d\log y_2$ to $-\lambda_1\wedge\lambda_2$.
There is no additional sign on the direct sum of strata.

\section{The full-subset normal complex}
\label{sec:normal-complex}

The smaller comparison requires a base projection.
Fix a product chart $U=\Delta^E\times Z$ and the projection $p:U\to Z$.
Here $\Delta$ is a complex disc and $m=|E|$.
Put $M^\bullet=\Gamma(Z,\Omega_Z^\bullet)$.
For $S\subseteq E$, the product stratum $U_S$ has projection $p_S:U_S\to Z$.
For an ordered set $F=E\setminus S$, define
\[
L(F)=\bigwedge^{|F|}\mathbb C^F[|F|],
\qquad
\omega_F=e_{f_1}\wedge\cdots\wedge e_{f_r},\qquad r=|F|.
\]
The line $L(F)$ lies in degree $-|F|$.
Deleting its $i$th basis vector gives
\[
\iota_{f_i}\omega_F=(-1)^{i-1}\omega_{F\setminus\{f_i\}}.
\]
Place this line before the coefficient complex and take the complete direct sum
\begin{equation}
\label{eq:normal-total}
C_E(M)=\bigoplus_{S\subseteq E}L(E\setminus S)\otimes M^\bullet,
\qquad A_E(M)=C_E(M)[-2m].
\end{equation}
On its $S$-summand, define
\begin{equation}
\label{eq:normal-differential}
D_A(\omega_F\otimes\beta)
=(-1)^r\omega_F\otimes d_Z\beta
+\sum_{e\in F}\iota_e\omega_F\otimes\beta.
\end{equation}
The residue term indexed by $e$ is in the $S\cup\{e\}$-summand.
The even shift in~\eqref{eq:normal-total} does not change this differential.

\begin{lemma}[The normal embedding]
\label{lem:normal-embedding}
The formula~\eqref{eq:normal-differential} defines a differential.
For adapted normal coordinates $x$, the map
\begin{equation}
\label{eq:normal-embedding}
j_x:A_E(M)\longrightarrow\mathcal R_Y(U),
\qquad
\omega_F\otimes\beta\longmapsto
\lambda_F^x\wedge p_S^*\beta
\end{equation}
is an injective chain map, where $\lambda_F^x=\bigwedge_{e\in F}d\log x_e$ in the given order.
Its image $N_x$ is a subcomplex of $\mathcal R_Y(U)$.
\end{lemma}

\begin{proof}
Two distinct deletions have opposite exterior signs.
The internal differential has signs $(-1)^r$ before a deletion and $(-1)^{r-1}$ after it.
The coefficient identity maps commute with $d_Z$.
These facts prove $D_A^2=0$.
The source degree of $\omega_F\otimes\beta$ is $|\beta|-r+2m$.
Since $m=r+|S|$, this equals the target degree $|\beta|+r+2|S|$.

Every normal logarithmic form is closed, so
$d(\lambda_F^x\wedge p_S^*\beta)=(-1)^r\lambda_F^x\wedge p_S^*(d_Z\beta)$.
Its $e_i$-residue is
$(-1)^{i-1}\lambda_{F\setminus\{e_i\}}^x\wedge p_{S\cup\{e_i\}}^*\beta$.
These are exactly the two parts of~\eqref{eq:normal-differential}.
They stay in the full direct sum.
Finally, the adapted coordinate exterior basis recovers $p_S^*\beta$ uniquely from its coefficient of $\lambda_F^x$.
Pullback along $p_S$ is injective on these forms, as restriction to a constant normal section recovers $\beta$.
Thus $j_x$ is injective.
\end{proof}

We next allow unit changes~\eqref{eq:adapted-change} that preserve the specified projection $p$.
The units may depend on both normal and base coordinates.
After restriction to a common chart domain, $(y_e,z_a)$ must still be adapted coordinates.
If a base chart changes by a biholomorphism depending only on the base coordinates,
its ordinary pullback on $M$ can be included in all formulas below.
A tangential coordinate change depending on the normal coordinates is outside this smaller comparison.
Proposition~\ref{prop:strict-residue-descent} still applies to the full residue complex in that case.

\begin{proposition}[Abstract transport and its geometric limit]
\label{prop:normal-transport}
A permutation $\pi$ induces a chain isomorphism
\[
P_\pi:A_{E_y}(M)\longrightarrow A_{E_x}(M),\qquad
(e_{f_1}\wedge\cdots\wedge e_{f_r})\otimes\beta
\longmapsto(e_{\pi(f_1)}\wedge\cdots\wedge e_{\pi(f_r)})\otimes\beta.
\]
It maps the $S$-summand to the $\pi(S)$-summand.
In ordered bases, its sign is the sign of sorting the list $\pi(f_1),\ldots,\pi(f_r)$.
These transports compose strictly.
The resulting abstract maps between normal images are
\[
T_{y\leftarrow x}=j_yP_\pi^{-1}j_x^{-1}:N_x\longrightarrow N_y.
\]
They also compose strictly.
In general they are not restrictions of geometric pullback on logarithmic forms.
\end{proposition}

\begin{proof}
Exterior contraction is natural under relabeling:
$\iota_{\pi(e)}P_\pi=P_\pi\iota_e$ on the appropriate summands.
The internal sign depends only on $|F|$, which is unchanged.
This proves the chain identity.
Functoriality of exterior powers gives $P_\rho P_\pi=P_{\rho\pi}$.
The formula for $T$ then cancels adjacent inverse maps in every composition.
The change~\eqref{eq:two-normal-unit} shows that geometric pullback does not preserve $N_x$.
Thus these abstract transitions need a further comparison with the ambient geometric maps.
\end{proof}

\section{A coherent comparison without choosing logarithms}
\label{sec:unfiltered-comparison}

A chain homotopy between chain maps $f,g:A\to B$ is a degree-minus-one map $H$ satisfying
$g-f=d_BH+Hd_A$.
It witnesses equality of the induced maps on cohomology.
The following contraction will supply such homotopies with an exact composition rule.
Its existence also determines the limited cohomological information carried by $A_E(M)$.

Identify the direct sum of lines $L(F)$ with the exterior algebra on the edge symbols,
placing each symbol in degree $-1$.
Let $\epsilon$ be the linear functional with $\epsilon(e)=1$ for every edge.
The edge part of $D_A$ is contraction $\iota_\epsilon$.
For $m>0$, set
\begin{equation}
\label{eq:average-contraction}
v_E=\frac1m\sum_{e\in E}e,\qquad
h(\omega_F\otimes\beta)=v_E\wedge\omega_F\otimes\beta.
\end{equation}
A wedge containing a repeated edge is zero.
Every nonzero term of $h$ adds an edge to $F$ and removes it from $S$.
Thus $h$ has degree $-1$, including after the shift $[-2m]$.

\begin{lemma}[Permutation-invariant contraction]
\label{lem:normal-contractible}
The identities $D_Ah+hD_A=\mathrm{id}$ and $h^2=0$ hold.
Moreover $P_\pi h=hP_\pi$.
Consequently $A_E(M)$ has zero cohomology for $m>0$.
\end{lemma}

\begin{proof}
For a vector $v$ and exterior form $\omega$, deleting each factor in $v\wedge\omega$ gives
$\iota_\epsilon(v\wedge\omega)=\epsilon(v)\omega-v\wedge\iota_\epsilon\omega$.
Here $\epsilon(v_E)=1$.
The internal terms cancel because the wedge degree changes from $r$ to $r+1$ under $h$.
This proves $D_Ah+hD_A=\mathrm{id}$.
The equality $v_E\wedge v_E=0$ proves $h^2=0$.
Every permutation fixes the average $v_E$, proving equivariance.
If $D_Aa=0$, then $a=D_A(ha)$, so every cycle is a boundary.
\end{proof}

On an overlap, use geometric pullback to place all forms in the same residue complex $B$.
For each chart $i$, let $A_i$ be its normal complex and $j_i:A_i\to B$ its embedding.
Let $P_{ij}:A_j\to A_i$ be the determinant transport.
It includes a base pullback if such a base change is present.
Assume these underlying atlas transitions satisfy $P_{ik}=P_{ij}P_{jk}$.
This is automatic for the stated coordinate permutations and compatible base charts.

\begin{theorem}[Overlap homotopies and their cocycle]
\label{thm:strict-homotopy-cocycle}
For $m>0$, define
\begin{equation}
\label{eq:overlap-homotopy}
H_{ij}=(j_j-j_iP_{ij})h_j:A_j^q\longrightarrow B^{q-1}.
\end{equation}
Then
\begin{align}
\label{eq:overlap-identities}
DH_{ij}+H_{ij}D_A&=j_j-j_iP_{ij},\\
H_{ik}&=H_{jk}+H_{ij}P_{jk}.
\end{align}
The construction uses $u_e^{-1}du_e$ but no choice of a logarithm of $u_e$.
The second identity is exact, so its triple-overlap defect is zero.
\end{theorem}

\begin{proof}
Write $\Delta_{ij}=j_j-j_iP_{ij}$.
Both terms are chain maps, hence $D\Delta_{ij}=\Delta_{ij}D_A$.
The first identity follows from
$D\Delta_{ij}h_j+\Delta_{ij}h_jD_A
=\Delta_{ij}(D_Ah_j+h_jD_A)=\Delta_{ij}$.
For the second, use $h_jP_{jk}=P_{jk}h_k$:
\[
H_{jk}+H_{ij}P_{jk}
=(j_k-j_jP_{jk})h_k
+(j_jP_{jk}-j_iP_{ij}P_{jk})h_k
=H_{ik}.
\]
The difference $\Delta_{ij}$ is a finite exterior expansion of~\eqref{eq:unit-differential}.
All its coefficients are globally defined regular forms on the overlap.
Thus no primitive of $d\log u_e$ enters the formula.
\end{proof}

For two aligned normal labels, let $\alpha_e=d\log u_e$ and
\[
W=(\lambda_1+\alpha_1)\wedge(\lambda_2+\alpha_2)
-\lambda_1\wedge\lambda_2.
\]
Suppressing a common base coefficient, the homotopy is
\begin{align*}
H(\omega_{\{1,2\}})&=0,\\
H(e_2\text{ on }S=\{1\})&=\tfrac12W\text{ on }S=\varnothing,\\
H(e_1\text{ on }S=\{2\})&=-\tfrac12W\text{ on }S=\varnothing,\\
H(1\text{ on }S=\{1,2\})
&=\tfrac12\alpha_1|_{D_2}+\tfrac12\alpha_2|_{D_1}.
\end{align*}
The last two terms belong to different summands.
Since $r_1W=\alpha_2|_{D_1}$ and $r_2W=-\alpha_1|_{D_2}$,
these formulas give the first identity in~\eqref{eq:overlap-identities} directly.
They also show that $H$ can decrease stratum depth by one.

\section{A chain map on a finite cover}
\label{sec:cech-map}

The overlap identities give an actual map on a complex of local sections.
Let $p:Y\to Z$ have a finite base cover $V_i$ with the following projected normal charts.
Each $p^{-1}V_i$ is identified over $V_i$ with an open subset of $\mathbb C^m\times V_i$
containing $\{0\}\times V_i$, and its normal divisors are the coordinate hyperplanes.
The map $p$ becomes projection to $V_i$, and the fixed number of normal directions is $m>0$.
The formulas above restrict to these open subsets; the zero section still recovers every base coefficient.
Normal labels may permute on overlaps, with their permutations satisfying the atlas cocycle.
The sheaves $A_i$ on $V_i$ consist of~\eqref{eq:normal-total}, with holomorphic forms on each base open set.
The maps $P_{ij}$ glue them to a sheaf of complexes $A$.
Explicitly, a section is a family of local sections with $a_i=P_{ij}a_j$ on overlaps.
Let $B=p_*\mathcal R_Y$.
Thus $B(V)=\mathcal R_Y(p^{-1}V)$.
After the identifications defining $A$, write each $j_i:A|_{V_i}\to B|_{V_i}$ and
$H_{ij}=(j_j-j_i)h$ on overlaps.

For a sheaf of complexes $F$, its alternating \v Cech cochains are
\[
\check C^p(F)^q=\prod_{i_0<\cdots<i_p}F^q(V_{i_0}\cap\cdots\cap V_{i_p}).
\]
The \v Cech differential is the alternating restriction sum
$(\delta a)_{i_0\ldots i_{p+1}}=
\sum_{r=0}^{p+1}(-1)^r a_{i_0\ldots\widehat i_r\ldots i_{p+1}}$.
Restrictions to the common overlap are implicit in this formula only.
The total complex has degree $p+q$ and differential
$\mathscr D a=\delta a+(-1)^p d_Fa$ on \v Cech degree $p$.
The equations $\delta^2=d_F^2=0$ and $\delta d_F=d_F\delta$ prove $\mathscr D^2=0$.

\begin{proposition}[The total comparison]
\label{prop:cech-comparison}
For $a\in\check C^p(A)^q$, define $\mathcal J a$ by its two components:
\begin{align}
\label{eq:cech-comparison}
(\mathcal J a)^{p}_{i_0\ldots i_p}
&=j_{i_0}a_{i_0\ldots i_p},\\
(\mathcal J a)^{p+1}_{i_0\ldots i_{p+1}}
&=(-1)^p H_{i_0i_1}a_{i_1\ldots i_{p+1}}.
\end{align}
This is a degree-zero chain map
$\operatorname{Tot}\check C(A)\to\operatorname{Tot}\check C(B)$.
No acyclicity condition on the cover is needed for this chain-map assertion.
\end{proposition}

\begin{proof}
The first component preserves both degrees.
The second raises the \v Cech degree by one and lowers the internal degree by one.
At unchanged \v Cech degree, the chain identity follows from $Dj_i=j_iD_A$.
At \v Cech degree $p+1$, the difference between restriction of the first component and application after $\delta$
is $(j_{i_1}-j_{i_0})a_{i_1\ldots i_{p+1}}$.
The internal differential of the second component has coefficient $-DH_{i_0i_1}$.
The second component applied to the internal differential has coefficient $H_{i_0i_1}D_A$ on the other side.
Thus the difference is
$(j_{i_1}-j_{i_0}-DH_{i_0i_1}-H_{i_0i_1}D_A)a_{i_1\ldots i_{p+1}}=0$.
At \v Cech degree $p+2$, all terms with later indices cancel as alternating restrictions.
The remaining coefficient is
$H_{i_1i_2}-H_{i_0i_2}+H_{i_0i_1}$, which vanishes by~\eqref{eq:overlap-identities}.
These exhaust the possible degrees.
\end{proof}

This construction gives a comparison into the global residue carrier.
It gives no quasi-isomorphism, meaning a map inducing isomorphisms on cohomology, with that carrier.
For $m>0$, the source sheaf $A$ is contractible by $h$.
Its \v Cech total complex is also contractible: on \v Cech degree $p$, use $(-1)^p h$.
The restriction terms cancel because this sign changes when $p$ increases.
The internal terms sum to the identity by Lemma~\ref{lem:normal-contractible}.
In contrast, on a connected nonempty chart $U$, the constant function $1$ is a nonzero class in $H^0(\mathcal R_Y(U))$.
It is closed and cannot be a boundary because the total complex has no negative degrees.
Consequently the normal comparison cannot recover even this degree-zero class.
When $m=0$, the product chart is $Z$ itself, $A=M$, and $j$ is the identity on its de Rham complex.
The contraction argument is unnecessary in that case.

\section{Homotopies that preserve stratum depth}
\label{sec:filtered-homotopies}

The preceding contraction adds a remaining edge and therefore removes a contracted edge.
This matters when the number of collisions is part of the structure.
Define the decreasing filtration
\[
F^aA=\bigoplus_{|S|\geq a}A_S,
\qquad
F^aB=\bigoplus_{|S|\geq a}B_S.
\]
Both differentials preserve it, since residues increase $|S|$.
A filtered map or homotopy sends $F^aA$ into $F^aB$ for every $a$.
The homotopy~\eqref{eq:overlap-homotopy} need not satisfy this condition.

On a sufficiently small overlap, a unit $u$ has a holomorphic logarithm $\ell$ with $e^\ell=u$.
To see this near a point, write $u=u_0(1+w)$ with $|w|<1$ and use the convergent power series for $\log(1+w)$.
Then $d\ell=u^{-1}du$.
Two logarithms differ by a locally constant element of $2\pi i\mathbb Z$.
After aligning the component labels, suppose $y_e=u_ex_e$ and choose logarithms $\ell_e$.
Write $\alpha_e=d\ell_e$.

Introduce a real parameter $t\in[0,1]$ only as a polynomial variable with its differential $dt$.
The forms
\[
\Lambda_e(t)=\lambda_e+t\alpha_e+dt\,\ell_e
=\lambda_e+d(t\ell_e)
\]
are closed for the total exterior derivative on the parameter and chart variables.
Their residues are $r_f\Lambda_e=\delta_{fe}$, where $\delta_{fe}$ is one for $f=e$ and zero otherwise.
Their other terms are regular.
Therefore the wedge construction~\eqref{eq:normal-embedding}, using $\Lambda_e(t)$,
gives a chain map into the polynomial parameter extension of the residue complex.

For a parameter form $a(t)+dt\wedge b(t)$, define
$I(a+dt\wedge b)=\int_0^1 b(t)\,dt$.
The integral means coefficientwise integration of a polynomial, so it is a finite algebraic operation over $\mathbb C$.
Direct differentiation gives
$dI+Id=\operatorname{ev}_1-\operatorname{ev}_0$.
Residues satisfy $r(dt\wedge b)=-dt\wedge rb$, and hence $RI+IR=0$.
Thus the same identity holds with $D=d+R$.

\begin{proposition}[Logarithmic overlap homotopy]
\label{prop:logarithmic-homotopy}
The degree-minus-one map
\begin{equation}
\label{eq:logarithmic-homotopy}
K_{xy}(\omega_F\otimes\beta)
=\int_0^1\sum_{i=1}^{r}(-1)^{i-1}\ell_{f_i}
\bigwedge_{\substack{j=1\\j\ne i}}^r
(\lambda_{f_j}+t\alpha_{f_j})\wedge p_S^*\beta\,dt
\end{equation}
is a filtered chain homotopy from $j_x$ to $j_y$.
All coefficients and forms on the right are restricted to $U_S$.
The remaining wedge factors keep their original order.
\end{proposition}

\begin{proof}
The coefficient of $dt$ in $\bigwedge_{j=1}^r\Lambda_{f_j}(t)$
is the integrand in~\eqref{eq:logarithmic-homotopy}.
Moving that $dt$ to the first position gives $(-1)^{i-1}$.
The parameter integration identity applied to the chain map just constructed yields
$DK_{xy}+K_{xy}D_A=j_y-j_x$.
Every term stays on $U_S$, so it preserves stratum depth.
It has one fewer differential form than its input image under $j_x$, and hence has degree $-1$.
\end{proof}

For two labels, the top-stratum formula is
\begin{equation}
\label{eq:two-log-homotopy}
K_{xy}(e_1\wedge e_2)
=\ell_1\lambda_2-\ell_2\lambda_1
+\tfrac12(\ell_1\alpha_2-\ell_2\alpha_1).
\end{equation}
On the $S=\{1\}$ summand it sends $e_2$ to $\ell_2|_{D_1}$.
On the $S=\{2\}$ summand it sends $e_1$ to $\ell_1|_{D_2}$.
It vanishes on $S=\{1,2\}$.
These formulas include their common base coefficient by wedging it on the right.

\begin{example}[A double overlap with no filtered homotopy]
\label{ex:period-obstruction}
Take $U=\Delta^2\times\mathbb C^*$ with base coordinate $z$.
Let $y_1=zx_1$ and $y_2=x_2$.
These functions and $z$ form an adapted coordinate system onto their open image.
The maps $j_x,j_y$ are defined on the same projected overlap.
There is no filtered chain homotopy $K$ from $j_x$ to $j_y$ on all of $U$.

To prove this, apply the homotopy equation to the coefficient-one element $e_1$ in $S=\{2\}$.
Its total degree is three, and $D_Ae_1=1$ in $S=\{1,2\}$.
A filtered degree-minus-one map sends that last element to zero:
its output would have degree three and depth at least two, whose minimum total degree is four.
The component of $Ke_1$ on $D_2$ is therefore a holomorphic function $f$ with
$df=dz/z$.
Restricting to $x_1=0$ gives the same equation on $\mathbb C^*$.
Along $z=e^{it}$, the integral of $df$ is zero because $f(e^{it})$ has the same endpoints.
The integral of $dz/z=i\,dt$ is $2\pi i$.
This contradiction proves nonexistence.
The unfiltered homotopy of Theorem~\ref{thm:strict-homotopy-cocycle} still exists on this overlap.
\end{example}

\section{The triple-overlap obstruction for two normal directions}
\label{sec:triple-obstruction}

Take three projected normal charts with labels aligned on a connected triple overlap.
Write their normal forms as
$\lambda^{(0)}=\lambda$,
$\lambda^{(1)}=\lambda+a$,
and $\lambda^{(2)}=\lambda+a+b$.
Choose logarithms $A_e$ for the change $0\to1$ and $B_e$ for $1\to2$.
Thus $a_e=dA_e$ and $b_e=dB_e$.
Let $C_e$ be the chosen logarithm for $0\to2$.
The coordinate cocycle gives
\begin{equation}
\label{eq:logarithm-obstruction}
c_e=A_e+B_e-C_e\in2\pi i\mathbb Z.
\end{equation}
The number $c_e$ is constant on the connected overlap.
Define the homotopy defect
$T=K_{01}+K_{12}-K_{02}$.
It has degree $-1$, and $DT+TD_A=0$ because the differences of the three embeddings cancel.
For a map $L$ of degree $-2$, the equation making it a secondary homotopy is
$DL-LD_A=T$.
Its minus sign is part of the cohomological convention.

\begin{theorem}[Exact filtered triple criterion]
\label{thm:filtered-triple-criterion}
Suppose both normal components meet on the triple overlap.
For the logarithmic homotopies~\eqref{eq:logarithmic-homotopy},
a filtered secondary homotopy exists if and only if $c_1=c_2=0$.
When these constants vanish, one can take
\begin{equation}
\label{eq:secondary-log-homotopy}
L(e_1\wedge e_2\otimes\beta)
=Q\,p_\varnothing^*\beta,
\qquad
Q=\tfrac12(A_2B_1-A_1B_2),
\end{equation}
on the ambient stratum and set $L=0$ on all other summands.
\end{theorem}

\begin{proof}
First assume $C=A+B$.
Substitution into~\eqref{eq:two-log-homotopy} gives
\begin{align*}
T(e_1\wedge e_2)
&=\tfrac12(B_1a_2-B_2a_1-A_1b_2+A_2b_1)\\
&=d\bigl(\tfrac12(A_2B_1-A_1B_2)\bigr)=dQ.
\end{align*}
On each one-edge summand, the three logarithms add to zero in $T$.
Thus $T$ vanishes there and on the empty-wedge summand.
For a base form $\beta$, differentiation of $Q\beta$ gives $dQ\wedge\beta+Qd\beta$.
The term $LD_A$ contributes $Qd\beta$ because the top remaining-edge count is two.
All edge terms of $LD_A$ vanish by the definition of $L$.
The residue of the regular form $Q\beta$ is zero.
Consequently $DL-LD_A=T$ on every summand.
The map stays on the ambient stratum and is filtered.

For arbitrary $C$, the restriction of $T$ to $e_1$ on $S=\{2\}$ is the constant function $c_1$ on $D_2$.
Its restriction to $e_2$ on $S=\{1\}$ is $c_2$ on $D_1$.
Let $L$ be any filtered map of degree $-2$.
On either of these degree-three inputs, $L$ must vanish, since a depth-one output has minimum total degree two.
The requested output degree would be one.
Their differentials are the degree-four unit at depth two.
The value of $L$ there must also vanish, since the requested degree two is below the depth-two minimum of four.
Hence $(DL-LD_A)$ vanishes on both one-edge inputs.
It cannot equal a nonzero $c_1$ or $c_2$ there.
This proves necessity and completes the criterion.
\end{proof}

This obstruction belongs to the chosen logarithmic homotopies.
Changing a logarithm on a double overlap by $2\pi i n_{ij,e}$ changes
$c_{ijk,e}$ by
$2\pi i(n_{ij,e}+n_{jk,e}-n_{ik,e})$.
Thus the obstruction to choosing logarithms with zero triple defect is the class of this constant two-cochain modulo those changes.
Here a two-cochain is a family on triple overlaps,
and two such families define the same class when their difference is the alternating change of a family on double overlaps.
On quadruple overlaps its alternating sum is zero, since each logarithm occurs twice with opposite signs.
If arbitrary holomorphic primitives of $d\log u_e$ are allowed in place of logarithms,
the permitted constants are in $\mathbb C$, and the corresponding class uses that constant coefficient space.
Component permutations transport these constants with their labels.

When additive logarithms have been chosen, the secondary homotopies also satisfy the next coherence identity.
For three successive logarithm vectors $A,B,C$, the scalar function in~\eqref{eq:secondary-log-homotopy} is bilinear:
\[
Q(A,B)+Q(A+B,C)=Q(B,C)+Q(A,B+C).
\]
Both sides expand to $Q(A,B)+Q(A,C)+Q(B,C)$.
Therefore the alternating quadruple-overlap sum of $L$ is zero.
Together, the embeddings, the homotopies $K$, and the secondary homotopies $L$
give a filtered map on the finite-cover total complexes of Section~\ref{sec:cech-map}.
For $a\in\check C^p(A)^q$, its components in \v Cech degrees $p,p+1,p+2$ are
\begin{align*}
(\mathcal J_{\log}a)^p_{i_0\ldots i_p}
&=j_{i_0}a_{i_0\ldots i_p},\\
(\mathcal J_{\log}a)^{p+1}_{i_0\ldots i_{p+1}}
&=(-1)^pK_{i_0i_1}a_{i_1\ldots i_{p+1}},\\
(\mathcal J_{\log}a)^{p+2}_{i_0\ldots i_{p+2}}
&=-L_{i_0i_1i_2}a_{i_2\ldots i_{p+2}}.
\end{align*}
The third term has internal degree $q-2$, so all three terms have total degree $p+q$.
To check the chain identity, use the calculation of Proposition~\ref{prop:cech-comparison}.
The new \v Cech-degree-$p+2$ defect is
$(-1)^p(\delta K-DL+LD_A)a$, which is zero.
At degree $p+3$ it is a signed copy of $(\delta L)a$, also zero by bilinearity of $Q$.
There are no components in higher degrees.
Thus no further homotopies are needed for this map.
The identities here are $\delta j=DK+KD_A$, $\delta K=DL-LD_A$, and $\delta L=0$;
they also express the total cocycle $j+K-L$ in the \v Cech complex of graded maps.

The constants do not obstruct the unfiltered construction.
Indeed, if $T$ is any closed degree-minus-one map out of $A$, then $-Th$ has degree $-2$ and
\[
D(-Th)-(-Th)D_A=T.
\]
To check this, use $DT=-TD_A$ and $D_Ah+hD_A=1$.
This primitive can decrease stratum depth, which is exactly the property excluded in Theorem~\ref{thm:filtered-triple-criterion}.

\section{What descends}
\label{sec:consumer-map}

The full residue complex glues by canonical pullbacks.
On a supplied projected atlas, the normal complex has determinant transports and the comparison~\eqref{eq:cech-comparison},
with the shift $[-2m]$ and every stratum retained.
For $m>0$, its contractibility excludes a quasi-isomorphism with the full residue complex.
When the homotopies must preserve stratum depth, the logarithmic construction instead gives the two-normal comparison
subject to~\eqref{eq:logarithm-obstruction}.

A construction with chiral coefficient states must supply those states on every stratum and the maps attached to coordinate changes.
Those maps must commute with residues and the internal coefficient differential on their actual domains.
The contraction $h$ used here also depends on the identity coefficient maps between all subsets.
For varying state spaces, adding an edge would require a map in the reverse stratum direction;
the formula $H=\Delta h$ requires such a compatible contraction to be constructed separately.
The logarithmic formula requires a differential-form action and compatible coefficient restrictions.
Changing the base projection through normal-dependent tangential coordinates also requires an additional coefficient comparison.
No such state or base-projection comparison is provided by the normal orientation line itself.
